\section{Discussion}
\label{sec:discussion}

\subsection{What Changed Relative to the Submitted Version}
The submitted version reported 96.88\,\% accuracy and 99.53\,\% recall for E-DDQN and a large gain over a standard DQN. Section~\ref{sec:leak} shows that numbers of that size follow from the data pipeline---oversampling before a row-level split, so that duplicated and temporally adjacent samples of the same sessions appear in training and test---and are reproduced by a generic classifier without reinforcement learning. Under a participant-conservative, nested and frozen protocol, the same data support only modest discrimination (best ROC-AUC \RFCostRocAucMean--\HGBCostRocAucMean), the corrected E-DDQN does not outperform the standard DQN, and the best cost-sensitive decision rule comes from a random forest. We have therefore replaced every numerical claim of the submitted version and withdrawn the claim that E-DDQN improves accuracy, recall, F1 and MCC over a standard DQN.

\subsection{Is E-DDQN Reinforcement Learning?}
With logged data the alarm cannot influence the subsequent telemetry, so the decision problem is a contextual bandit (Section~\ref{sec:formulation}). E-DDQN remains a valid value-based learner of a cost-sensitive detection policy, and its episodes, transitions and terminal states are now defined consistently with the data. However, (i) the asymmetric reward is mathematically equivalent to a 5:1 cost-sensitive Bayes threshold, so it is a statement about the operating point rather than a different learning objective; (ii) bootstrapping cannot improve the optimal policy, and empirically strong bootstrapping ($\gamma$ close to 1) made the learned margins noisy and the validation-based selection over-optimistic; and (iii) under bandit feedback the agent observes the reward of one action only, although labels reveal both, which is why full-information replay (storing each counterfactual action with its own reward) is the principled replacement for teacher forcing. Adaptive sequential control would require an intervention in the loop---for example an alarm that triggers a rendering mitigation whose effect on later telemetry and symptoms is measured---which no public cybersickness dataset provides.

\subsection{The Safety Objective}
A 5:1 cost makes an alarm the cheaper decision whenever $P(y=1\mid s)\ge1/6$. With a prevalence of 0.44 and weakly informative inputs, almost every second satisfies this, so the cost-optimal policy of every model is close to ``always alarm'': recall is high, but \DOneFPRMin--\DOneFPRMax\,\% of non-event seconds also trigger an alarm (mean over seeds, rule D1). High recall alone is therefore not evidence of a useful safety monitor, and the comparison has to be made on cost, specificity at matched recall and threshold-free ranking. On those criteria the intended safety objective---near-complete detection with a controlled false-alarm rate---is \emph{not} achieved by any model on these data. The obstacle is the information content of game telemetry: once participants are not shared between training and test, exposure time, accumulated motion and the pre-session questionnaire carry most of the usable signal, and they distinguish individual susceptibility only coarsely.

\subsection{The Adversarial Threat Model}
The threat model is motivated by documented attacks on VR perception and frame pacing~\cite{odeleye2021framerate,tseng2022dark,kundu2025securing}, but UFFCSData contains no attacked sessions, so the binary target is a cybersickness label, not an attack label. What the data support is a stress test of the detector under bounded manipulation of the frame-rate, orientation and locomotion channels; field-of-view manipulation cannot be tested because the field of view never varies in the data. Under these tests, \AttAucFracSmall\,\% of model--scenario combinations change ROC-AUC by less than 0.03 and none by more than \AttAucMaxAbsChange, because the detectors rely mostly on exposure and questionnaire inputs, whereas an adaptive adversary who knows the labels can lower recall on event seconds by \AttEvasionDropMin--\AttEvasionDropMax{} percentage points, depending on the model. Demonstrating detection of real rendering attacks requires data in which manipulations are applied to participants under ethical approval and their physiological effects are measured.

\subsection{Limitations}
(i) A single public dataset with \nSessionsEval{} curated sessions and \nGroupsEval{} independent groups; confidence intervals are wide and fold-to-fold variation is large. (ii) Participant identity cannot be verified; linkage groups are conservative but cannot exclude a participant who answered the questionnaire differently in non-consecutive sessions. (iii) The label is a coarse self-report sampled at 1\,Hz and is nearly monotone within a session. (iv) Perturbations are synthetic and feature-space; their physiological effect is unknown. (v) Hyper-parameter search used an identical but modest budget (\HPOBudget{} configurations per family); a larger budget could change the ranking of the neural models, although the selection-optimism analysis suggests that additional search would mostly increase optimism. (vi) The October-2019 sessions were used only for development and could not serve as an external test set because their camera-orientation channels were not logged.

\subsection{Implications}
Accuracies reported on gameplay cybersickness datasets---including ours in the submitted version---should be read together with the split protocol: with 1-Hz samples and near-monotone labels, any split that shares sessions or participants between training and test inflates performance to near-perfect levels. For safety-prioritised detection, the asymmetric reward should be treated as what it is, a cost-sensitive operating point, and RL detectors should be compared with cost-weighted supervised models at matched recall or cost under participant-independent evaluation.
